% Section 1, part 3: inputs, what is not claimed, strategy, plan of the note.

%-------------------------------------------------------------------------------
\subsection{Inputs}
\label{sec:int-inputs}

The results are conditional on the following inputs. Statement numbers refer to the arXiv versions named in the bibliography.
\begin{description}
\item[(Q)] \emph{Quoted theorems.} Topological recursion, applied only to the model curve $x=\Wv(p)$, $y=p$ and its rescalings $x=\kappa\Wv$ (\cref{sec:mod}): \cite[Def.~4.2, (4-9), with the kernel (4-3); Theorems~4.2 and~4.6; Theorem~5.1, (5-14), with \S5.3, (5-27)--(5-30)]{Eynard:2007kz}; the theorem \cite[Theorem~4.1, with (4-1)--(4-3)]{DuninBarkowski:2012bw}, with its graph weights \cite[\S2.3, (2-22)--(2-26)]{DuninBarkowski:2012bw} (the dilaton leaves (2-23) and the factor $1/|\Aut|$ of (2-26)), as recalled in \cite[\S2.6, (2-30)--(2-32), (2-34), (2-35)]{Dunin-Barkowski:2015caa}, with matrices in the left-action convention of \cite[Remark~2.1]{Dunin-Barkowski:2015caa} (upper index the row), and with the action of \cite[\S2.3]{Dunin-Barkowski:2015caa} read with the factor $1/|\Aut\Gamma|$ and with the dilaton leaves of \cite[Def.~2.13]{PPZ15}, that is, as the action of \cite[Defs.~2.2, 2.13]{PPZ15}; it is applied to $x=4\Wv$, $y=p$, whose critical points $p^5=v$ have the pairwise distinct critical values $5p_c^4$. Its hypotheses, namely symplecticity, \cite[(2-32)]{Dunin-Barkowski:2015caa} and the edge condition \cite[(4-1)]{DuninBarkowski:2012bw}, are proved for this curve in \cref{prop:mod-Rtr}; the form \cite[(2-33)]{Dunin-Barkowski:2015caa} of the edge condition is misprinted (it has $+\sum_k$ in place of $\delta_{ij}-\sum_k$) and is not used, nor are \cite[Prop.~4.3]{Dunin-Barkowski:2015caa} and the sentence preceding \cite[Theorem~8.1]{Dunin-Barkowski:2015caa}; and \cite[Theorem~7.1: (7.6), (7.7) with its factor $1/\hbar^2$, (7.9), (7.10), with the kernel (1.1) and the assumptions of \S1.1]{ABDKS25}. The $R$-matrix action: \cite[\S0.2, (1); Def.~0.3, in particular (iii); (5); Remark~2.1 and (6); Def.~2.2, with the convention stated after it (p.~20); Def.~2.5, with the convention stated after it; Prop.~2.12, with the formula (8) in its proof; Def.~2.13; Prop.~2.14]{PPZ15}; Props.~2.12 and~2.14 are used as statements over $\bbQ$ and transferred to the other coefficient rings in \cref{def:set-graphsum}; (8) enters only through \cref{lem:set-pushins}.
\item[(QRR)] \emph{The $R$-matrix along the line: flatness and the value at the large-radius point} (\cref{thm:set-qrr}). The $R$-matrix $R$ of (T) has coefficients in $\End(\bbH)\otimes\bbQ(\zeta)\fps q$ and satisfies: (\emph{flatness}) for every colour $i$ the column $R\,\tilde e_i\,\re^{u_i/z}$ solves the quantum differential equation $z\partial_Tv=H*v$ along the line (\cref{def:set-rows}~(b)); (\emph{value}) $R(q=0)=\diag(h_{\bw_0},\dots,h_{\bw_4})$ in the normalised canonical basis at $q=0$. Flatness is a theorem about the $R$-matrix of the cohomological field theory, not a convention; the convention (Givental's) is only that the legs of the graph sum carry $R^{-1}(\psi)$. Flatness is used in \cref{prop:set-cyclicLR}~(iii) and in the last sentence of \cref{cor:set-cyclicPF5}~(b), on which \cref{lem:gau-colrows}, \cref{thm:gau-R0}~(iii), \cref{cor:gau-rowsNN}, \cref{thm:con-glue} and \cref{lem:mai-Psi} rest. The value is used for every colour, in the form $\Phi_i(z,1)=h_{\bw_i}(z)$ (\cref{eq:set-qrr}): in \cref{lem:gau-colrows} and, for colour $0$, in the glue \cref{thm:con-glue}. Sources for flatness: \cite[Lemma~5.2, (16)]{GJR18}, which is \cref{prop:set-cyclicLR}~(iii) at $\lambda=1$; \cite[\S4.2, (12)--(13)]{GJR17}; in general \cite[Section~1, Proposition (a)--(c)]{Giv01} and \cite[Props.~7.7, 8.2~(ii), 8.4, 8.7~(ii)]{MR2917177}. Sources for the value: \cite[(17), in the proof of Lemma~5.2]{GJR18}; \cite[(19)--(20), in the proof of Lemma~4.6, at $t=0$]{GJR17}, which state it for this theory; \cite[\S2.2, Theorem~1 and the Remark after it]{Giv01}; \cite[Theorem~1, Corollary~1]{MR2276766}, which give the tangent factor and the Euler-class twist separately.
\item[(M)] \emph{The genus-zero mirror theorem} \cite[Theorem~11.6]{MR1408320}\draftnote{numbers from arXiv v2; IMRN numbering not compared}, in the Calabi--Yau case $n=4$, $r=1$, $l_1=5$, $\lambda'=0$, with the set-up of \cite[\S9]{MR1408320} and \cite[Prop.~11.4]{MR1408320}: at $\lambda_i=\zeta^i$, $J(T(q)H,z)=I(q,z)/C_0(q)$ with $T=I_1/C_0$ (\cref{thm:set-mirror}).
\item[(G0)] \emph{The genus-zero axioms of the twisted theory}: the genus-zero twisted classes form a cohomological field theory with unit for $\eta$ \cite[\S4.1]{GJR18}, \cite[Def.~0.3]{PPZ15}, polynomial and $S_5$-symmetric in $\lambda$, with the string, divisor and splitting axioms, the genus-zero topological recursion relation and the quantum differential equation \cite[\S5.2, the display before (14), of which (14) is the form in the $q$-direction]{GJR18} (\cref{thm:set-G0}).
\item[(T)] \emph{Teleman's reconstruction, with the axioms of the twisted theory in all genera} (\cref{thm:set-GT}). (i) At $\lambda_i=\zeta^i$ the classes $\Omega_{g,n}$ of \cref{def:set-fq}, in all genera, form a cohomological field theory with flat unit $1$ for $\eta$ \cite[\S4.1]{GJR18}, \cite[Def.~0.3]{PPZ15}, and satisfy the divisor equation: for $2g-2+n>0$ the part of degree $d$ of $\pi_*\Omega_{g,n+1}(\gamma_1,\dots,\gamma_n,H)$ is $d$ times that of $\Omega_{g,n}(\gamma_1,\dots,\gamma_n)$, where $\pi$ forgets the last point; so the theory at $\Qn=\re^{T(q)}$ is the theory at the point $t=T(q)H$ with $\Qn=1$. (ii) This theory, which is semisimple (\cref{cor:set-cyclicPF5}~(a)), is the graph sum of \cref{def:set-graphsum}, $\Omega=R.\TFT$, for an $R$-matrix $R$ with coefficients in $\End(\bbH)\otimes\bbQ(\zeta)\fps q$, symplectic and with $R|_{z=0}=1$ \cite[Theorem~(2) in (1.15)]{MR2917177}, \cite[\S4.1, Theorem~4.1]{GJR18}; this is the $R$-matrix of \cref{def:set-rows}, to which (QRR) refers. In \cref{thm:int-main,thm:int-D} (T) identifies that $R$-matrix, whose rows define $T_\Gamma$ and $\Dn$; in \cref{cor:int-gap} it identifies the total with $F_g$.
\end{description}
In addition, the formula \cite[(1.7)]{AC96}, with $\kappa_0=2g-2+n$ \cite[p.~5]{AC96}, is used in part (i) of the dilaton lemma \cref{lem:set-dilaton}. The definitions use \cite[\S0.2, (5)]{LP18}, \cite[\S0.2, (5), (10)]{LP20} and \cite{MR1666787}; \cite{ZZ08} is cited for the series of \cref{lem:set-periods} only. The inputs are read with the following definitions and conventions of their sources: the normalisation of the $J$-function \cite[Section~6, (10)]{MR2276766} in (M) (\cref{thm:set-mirror}); the free energies \cite[Def.~4.3, (4-10)]{Eynard:2007kz} (\cref{conv:mod-kernel,cor:mod-trF}) and the set-up \cite[(5-1), the assumption before (5-4), (5-13)]{Eynard:2007kz} of \cite[Theorem~5.1]{Eynard:2007kz} (\cref{thm:mod-var}); the recursion \cite[(2-26)--(2-29)]{Dunin-Barkowski:2015caa}, to which the recollection of \cite[Theorem~4.1]{DuninBarkowski:2012bw} in \cite[\S2.6]{Dunin-Barkowski:2015caa} refers (\cref{conv:mod-kernel}); and the reading of the integrals \cite[(2-31), (2-32)]{Dunin-Barkowski:2015caa} as formal Laplace series (\cref{lem:mod-laplace,prop:mod-Rtr}): termwise Gaussian integration of the holomorphic part, as in \cite[(3-57)]{DuninBarkowski:2012bw}, and the value $\mathsf G_{-2}=-\vartheta$ for the polar term $w^{-2}\dd w$ (cf.\ \cite[(3-64)]{DuninBarkowski:2012bw}); and the edge condition \cite[(4-1)]{DuninBarkowski:2012bw} read in the coordinates $w_i$ of \cite[\S2.6]{Dunin-Barkowski:2015caa}, that is, as \eqref{eq:mod-edge} (\cref{prop:mod-Rtr}~(iii)).

\emph{Where the inputs enter.} (M) and (G0) enter only in \cref{sec:set} and its appendix \cref{sec:set-lrproof} (\cref{thm:set-mirror,prop:set-cyclicLR,cor:set-cyclicPF5}); (QRR) at the places named above: flatness in \cref{prop:set-cyclicLR}~(iii) and \cref{cor:set-cyclicPF5}~(b), the value in \cref{lem:gau-colrows} and \cref{thm:con-glue}; (T) in \cref{thm:set-GT} and \cref{cor:mai-gap}; topological recursion only in \cref{sec:mod}; the details are in \cref{sec:mai-thm}. \Cref{sec:gau} uses no quoted input beyond (M), (G0) and (QRR), which enter through the rows of \cref{sec:set} (flatness, \cref{def:set-rows}~(b), in \cref{thm:gau-R0}~(iii)--(iv)) and in \cref{lem:gau-colrows} (the value; hence in \cref{cor:gau-rowsNN}~(iii) and \cref{cor:gau-colours}); the transport lemmas of \cref{sec:tra} use \cite[Defs.~2.2, 2.13]{PPZ15} and the dilaton lemma, and Theorem~R adds no input; Theorem~P(a) (\cref{thm:con-Pa}) and part (a) of \cref{thm:int-gamma} use no quoted input.

\emph{Conditionality.} Theorem~A holds modulo the quoted theorems (Q) of \cite{Eynard:2007kz}, \cite{DuninBarkowski:2012bw} (as recalled in \cite{Dunin-Barkowski:2015caa}), \cite{ABDKS25} and \cite{PPZ15} with the statement numbers above, the flatness of the $R$-matrix and its value at the large-radius point (QRR), the genus-$0$ mirror theorem (M) and the genus-$0$ axioms (G0) of the twisted theory, Teleman's reconstruction with the axioms of the twisted theory in all genera (T), and the formula \cite[(1.7)]{AC96}; every other step is proved in the note.

In particular Theorem~U (over every field of characteristic $0$, \cref{thm:gau-U}) and every lemma of the chain are proved here; Stirling's series enters only through its formal proof as the stationary phase of Euler's integral (\cref{lem:fsp-stirling}), and the implicit function theorem only in its formal version (\cref{lem:fsp-ift}). \Cref{cor:int-gap} uses, beyond \cref{thm:int-main}, only (T), part of \cite{PPZ15} (listed after \cref{cor:mai-gap}) and proved statements. \Cref{thm:int-D} uses the $R$-matrix that (T) identifies only through $R|_{z=0}=1$, its coefficient ring $\bbQ(\zeta)\fps q$ and (QRR), that is, flatness (\cref{def:set-rows}~(b)) and the value at $q=0$, and the genus-zero data of \cref{sec:set}, which rest on (M) and (G0). \Cref{thm:int-gamma} uses no quoted input; the identity \cite[24.4.1]{DLMF} of Bernoulli polynomials used in part (b) is derived from their generating function in the proof of \cref{lem:fsp-dlmf5118}. No finite computation enters a proof; the computations that accompanied the work are summarised in \cref{sec:evi}.

All other citations are further sources, comparisons, conventions (\cref{sec:sgn}) or context, and enter no proof; among them \cite[Corollary~1.2]{Hock25} (\cref{rmk:mod-hock}), \cite[Theorem~4.4, (4-23); Corollary~4.1, (4-26)--(4-29); (A-1), (A-5), (A-25); Theorem~4.7, (4-30); (5-15)]{Eynard:2007kz}, \cite[(26)]{LP20}, \cite[Corollary~11.7]{MR1408320} (a misprinted form of (M)), \cite[Props.~6.1--6.2]{GR19}, \cite[\S2.4.3, (17); \S2.4.4, (19)--(20)]{LP20} and \cite{CGYZ26}.

%-------------------------------------------------------------------------------
\subsection{What is not claimed}
\label{sec:int-notclaimed}

\begin{enumerate}
\item Nothing is claimed about the quintic threefold or its graph sums; the formal quintic is a different theory. The conifold gap of the quintic threefold is \cite[Theorem~2]{CGYZ26}.
\item Nothing is claimed for frames other than those of \cref{def:set-frames} and the rescalings of their two periods (\cref{sec:sgn-lit}), for $g=1$, or away from $u=0$.
\item Nothing is claimed about the analytic continuation of the series $F_g(q)$ to a neighbourhood of $u=0$. The conifold expansion $\Cexp_{\alpha,\beta}$ is a formal change of frame on the ring $\Nring$ (\cref{def:set-confexp}), applied to the graph-sum representative $F_g^{\Nring}$ of $F_g$ (\cref{def:mai-Fcf}); this is the sense of ``conifold expansion'' and of ``conifold gap'' in \cref{cor:int-gap}.
\item The second key statement of \cite[Remark~1.9]{CGYZ26} is not addressed (the Model Lemma, \cref{def:red-ML}, concerns the model), and \cref{cor:int-gap} is not claimed to be literally the first, which is named there but not stated.
\item No novelty is claimed for the formal stationary-phase lemmas of \cref{sec:fsp}, proved for completeness, for the formula of $\Ginf$ \cite[Lemma~7.7, (37)]{GJR18}, or for the change of chart \cite[\S10]{MR1901075}.
\item Topological recursion is not applied to a mirror curve of the quintic or of the formal quintic, and \cite{FCCZ26} is not used.
\end{enumerate}

%-------------------------------------------------------------------------------
\subsection{Strategy}
\label{sec:int-strategy}

\emph{Reduction to the model (\cref{sec:gau,sec:tra,sec:red}).} In a conifold frame the rows satisfy a cyclic system with one simple pole at $u=0$ (\cref{prop:gau-quintic}), from \cref{prop:set-cyclicLR} and the cyclic form of $\PFfour$ (\cref{thm:gau-cyc}); the model pulled back by $v(u)$ has the same product of coefficients, which fixes $\tau$ (\cref{prop:gau-equalprod}). Theorem~U gives the regular gauge between them, and the gauge identity singles out the symplectic gauge $\Pc$ (\cref{thm:gau-G,prop:gau-symp}). Its edges lie in $\sKo\tensor\sKo$ for a slice $\sK$ defined by the vanishing solution of the model (\cref{prop:tra-can}), and Theorem~R (\cref{thm:red-main,cor:red-can}) reduces the regularity of the edge terms of the flat sum of $\Pc.\Omega^{\mathrm{mod}}_{v(u)}$, and the vertex law, to the Model Lemma on $\sK$ (\cref{def:red-ML}).

\emph{The model (\cref{sec:mod}).} The model cohomological field theory is topological recursion on its curve (\cref{thm:mod-tr}); insertions from $\sK$ become legs with pairings polynomial in $v$ (\cref{prop:mod-LD,thm:mod-V}), and the free energies are $c_hv^{2-2h}$ (\cref{thm:mod-gap}); this is the Model Lemma (\cref{thm:mod-MLK}).

\emph{The normalisation (\cref{sec:flh,sec:con,sec:gam}).} The comparison theorem (\cref{thm:flh-main}) identifies $T_\Gamma$ with the flat term, graph by graph, if $\Dn_0\equiv1\bmod z^{3g-2}$; Theorem~P (\cref{thm:con-Pa,thm:con-Pb}) and the glue (\cref{thm:con-glue}) turn $\Dn=1$ into part (a) of \cref{thm:int-gamma} (\cref{thm:gam-D}). \Cref{thm:mai-main} assembles the parts.

%-------------------------------------------------------------------------------
\subsection{Plan of the note}
\label{sec:int-plan}

\Cref{sec:set} defines the objects and states (M), (G0), (QRR) and (T); \cref{sec:gau} proves Theorem~U, the gauge identity and symplecticity; \cref{sec:tra} proves the transport lemmas and \cref{sec:red} Theorem~R; \cref{sec:mod} proves the Model Lemma; \cref{sec:flh} compares the flat and the outer sums; \cref{sec:con} proves Theorem~P and the glue, and \cref{sec:gam} \cref{thm:int-gamma,thm:int-D}. \Cref{sec:mai} proves \cref{thm:int-main,cor:int-gap}, and \cref{sec:cmp} records the natural factorisation and universality. The appendices prove the large-radius cyclic structure (\cref{sec:set-lrproof}), the cyclic form of $\PFfour$ (\cref{sec:gau-cycproof}), the $u$-adic convergence lemma (\cref{sec:tra-convproof}) and the lemmas of formal stationary phase (\cref{sec:fsp}), and record the sign conventions (\cref{sec:sgn}) and the finite computations (\cref{sec:evi}).


\subsection*{Acknowledgements}

This work was supported by the UKRI grant UKRI2394 and by the Leverhulme Trust grant LIP-2023-007.
